\documentclass[10pt,a4paper]{amsart}
\usepackage[margin=19mm]{geometry}
\usepackage{amsmath,amssymb,mathtools}
\usepackage[scr=rsfs,cal=euler]{mathalfa}
\usepackage{xfrac}
\usepackage[unicode,hidelinks,bookmarks=false]{hyperref}
\newcommand{\ie}{i.e.\ }
\newcommand{\cf}{cf.\ }
\newcommand{\resp}{respectively\ }
\newcommand{\Subsection}{\subsection}
\providecommand{\nobreakdash}{\nobreak\hbox{-}}
\setlength{\emergencystretch}{2em}
\multlinegap0pt
\usepackage{bm}
\usepackage[all]{xy}
\usepackage{tcolorbox}
\usepackage{tikz}
\usepackage{amssymb}
\usepackage{multirow}
\usepackage{float}
\usepackage{xcolor}
\newtheorem{lemma}{Lemma}
\title{Polynomial algebra from the Lie algebra reduction chain $\mathfrak{su}(4) \supset \mathfrak{su}(2) \times \mathfrak{su}(2)$: The supermultiplet model}
\author{Rutwig Campoamor-Stursberg, Danilo Latini, Ian Marquette, Junze Zhang, Yao-Zhong Zhang}
\date{Source: arXiv:2503.04108v1 (CC BY 4.0)}
\begin{document}
\maketitle

\noindent\emph{Source and licence.} Polynomial algebra from the Lie algebra reduction chain $\mathfrak{su}(4) \supset \mathfrak{su}(2) \times \mathfrak{su}(2)$: The supermultiplet model. Version arXiv:2503.04108v1 (CC BY 4.0), distributed by arXiv under the Creative Commons Attribution 4.0 International licence, http://creativecommons.org/licenses/by/4.0/, as recorded in the arXiv licence metadata for this exact version and shown on the abstract page https://arxiv.org/abs/2503.04108v1. Full licence notice: \texttt{licenses/arxiv-2503.04108-CC-BY-4.0.txt}.\par\smallskip

\noindent\textit{Editorial scope.} Counted unit: the article's lemma 3.9 (source0.tex lines 998--1009). The article's complete original proof of it is delivered with it (source0.tex lines 1012--1039, one \texttt{proof} environment of 28 lines). No mathematics is written, repaired or re-derived: the statement and the proof are the article's own lines, in the article's own notation, and no retained line is substituted or re-flowed.  Retained uncounted as the source-local context the proof needs: lines 235--238; lines 239--242; lines 470--504; lines 984--988.  Nothing was omitted from any retained span, so the package carries no omission record.  The retained lines cite 1 works, whose entries are transcribed verbatim from the article's own bibliography source in the e-print and delivered with short editorial labels recorded in the item's reference mapping.  No retained line contains a figure environment, an image-inclusion command or drawing code, so no image is delivered and the package declares no asset dependency.  Editorial formatting only, in addition: an AMS \texttt{amsart} preamble that restates the article's own declarations and macros used by the retained lines, namely the environments \texttt{lemma} on separate counters, together with the standard packages those lines need.  The retained environments print in this document as Lemma 1 in order of appearance, whereas the article prints the same items with the numbering of its own full document.  The complete native source of the article as distributed in the arXiv e-print for this exact version is bundled unchanged as \texttt{sources/174307/source0.tex}.
\begin{equation}
	\{f,g\}=\sum_{i,j,k=1}^n C_{ij}^k x_k \partial_{x_i} f \partial_{x_j} g  \, ,
	\label{FPB}
\end{equation} where $\partial_{x_i} := \partial/\partial x_i.$ See, for instance, \cite{MR2906391}.  The algebra of smooth functions inherits the structure of a Poisson algebra that contains a subalgebra isomorphic to $\frak{g}$.  From this, one can define the linear Poisson bracket (a.k.a. Poisson-Lie structure) relations on $\mathfrak{g}^*$ as follows:

\begin{equation}
	\{x_i,x_j\}= \sum_{k=1}^nC_{ij}^k x_k  \, .
	\label{xcoord}
\end{equation} 

\subsection{The reduction chain $\mathfrak{su}(4) \supset \mathfrak{su}(2) \times \mathfrak{su}(2)$}
\label{subsec3.1}
The fifteen-dimensional Lie algebra $\mathfrak{su}(4)$ is given in terms of the basis: 
\begin{equation}
\{S_1, S_2, S_3, T_1, T_2, T_3, Q_{11}, Q_{12}, Q_{13}, Q_{21}, Q_{22}, Q_{23}, Q_{31}, Q_{32}, Q_{33}\} 
\label{gen}
\end{equation}
and the $\mathfrak{su}(2) \times \mathfrak{su}(2)$ subalgebra is spanned by the six generators $\{S_1, S_2, S_3, T_1, T_2, T_3\}$. The commutation relations, for $1 \leq i,j,k,\alpha,\beta \leq 3$, are the following:
\begin{align}
[S_i, S_j]&={\rm i} \epsilon_{ijk}S_k \qquad [T_\alpha, T_\beta]={\rm i} \epsilon_{\alpha \beta \gamma}T_\gamma \qquad [S_i, T_\alpha]=0 \label{a} \\
[S_i, Q_{j\alpha}]&={\rm i} \epsilon_{ijk} Q_{k \alpha} \quad \,\, [T_\alpha, Q_{i\beta}]={\rm i} \epsilon_{\alpha \beta \gamma} Q_{i \gamma} \quad [Q_{i \alpha}, Q_{j \beta}]=\frac{{\rm i}}{4}(\delta_{\alpha \beta} \epsilon_{ijk}S_k+\delta_{ij}\epsilon_{\alpha \beta \gamma}T_\gamma) \, , \label{b}
\end{align}
where summation over repeated indices is understood and $\rm i = \sqrt{-1}$. Here $\delta_{ij}$ is the Kronecker delta and $\epsilon_{ijk}$ is the Levi-Civita symbol. Under the basis $\eqref{gen},$ the Lie algebra $\mathfrak{su}(4)$ admits the (vector space) decomposition $ \mathfrak{su}(4) = \mathfrak{g}_1 \oplus \mathfrak{g}_2 \oplus \mathfrak{g}_3$ with 
\begin{align*}
   \mathfrak{g}_1 = \mathrm{span} \{S_1,S_2,S_3\} , \text{ } \mathfrak{g}_2 = \mathrm{span} \{T_1,T_2,T_3\} \text{ and } \mathfrak{g}_3 = \mathrm{span} \{Q_{11},\ldots,Q_{33}\}.
\end{align*} In this case, we consider the following linear coordinates: 
\begin{align}
\boldsymbol{x}:&=\{x_1, x_2, x_3, x_4, x_5, x_6,x_7,x_8,x_9,x_{10},x_{11},x_{12}, x_{13}, x_{14}, x_{15}\} \nonumber \\ &=\{s_1,s_2,s_3,t_1,t_2,t_3,q_{11},q_{12},q_{13},q_{21},q_{22},q_{23},q_{31},q_{32},q_{33}\} \, ,
\label{lincoord}
\end{align}
 and restrict to consider the commutant related to the subset of elements: 
 \begin{equation}
 \boldsymbol{x}':= \{x_1, x_2, x_3, x_4, x_5, x_6\} = \{s_1, s_2, s_3, t_1, t_2, t_3\} 
 \label{subalgebcoordin}
 \end{equation}
with respect to the Berezin bracket  \eqref{FPB}. In the commutative setting, the brackets of coordinates \eqref{xcoord} explicitly read:
\begin{equation}
\begin{split}
\{s_i, s_j\}&={\rm i} \epsilon_{ijk}s_k \qquad \{t_\alpha, t_\beta\}={\rm i} \epsilon_{\alpha \beta \gamma}t_\gamma \qquad \{s_i, t_\alpha\}=0 \\
\{s_i,q_{j\alpha}\}&={\rm i} \epsilon_{ijk} q_{k \alpha} \quad \,\, \{t_\alpha, q_{i\beta}\}={\rm i} \epsilon_{\alpha \beta \gamma} q_{i \gamma} \quad \{q_{i \alpha}, q_{j \beta}\}=\frac{{\rm i}}{4}(\delta_{\alpha \beta} \epsilon_{ijk}s_k+\delta_{ij}\epsilon_{\alpha \beta \gamma}t_\gamma) \, .
\label{classicalrels}
\end{split}
\end{equation}
 These relations will serve as the building blocks for deriving the polynomial algebra relations.
\vskip 0.5cm

Assume that $\mathfrak{g} = \mathfrak{g}_1 \oplus \mathfrak{g}_2 \oplus \mathfrak{g}_3$ is a semisimple Lie algebra  with  \begin{align}
\nonumber
    [\mathfrak{g}_1,\mathfrak{g}_1] & \,\subset \mathfrak{g}_1, \quad  [\mathfrak{g}_2,\mathfrak{g}_2] \subset \mathfrak{g}_2, \quad  [\mathfrak{g}_1,\mathfrak{g}_2] = \{0\},  \\
    [\mathfrak{g}_1,\mathfrak{g}_3] & \,\subset \mathfrak{g}_3, \quad  [\mathfrak{g}_2,\mathfrak{g}_3] \subset \mathfrak{g}_3, \quad  [\mathfrak{g}_3,\mathfrak{g}_3] \subset \mathfrak{g}_1 \oplus \mathfrak{g}_2. \label{eq:c1}
\end{align} 

\begin{lemma}
\label{3.9}
  Let $\mathfrak{g}^*$ be its dual admitting the same relations as in $\eqref{eq:c1}$, and let $\mathcal{Q}_\mathfrak{g}(d) \subset S(\mathfrak{g})^{\mathfrak{g}_1 \oplus \mathfrak{g}_2 }.$ For any non-zero indecomposable monomials $p \in \mathcal{Q}_{i_1 + i_2  + i_3}$ and $q \in \mathcal{Q}_{i_1' + i_2' + i_3'}$,   \begin{align}
  \nonumber
   \mathcal{G} \left(\{p,q\}\right) =     (i_1+i_1'-1,i_2+i_2',i_3+i_3') \tilde{+}(i_1+i_1',i_2+i_2'-1,i_3+i_3') \tilde{+}    \left\{\begin{matrix}
(i_1+i_1'+1,i_2+i_2',i_3+i_3'-2)  & \text{ if } \text{ }  [\mathfrak{g}_3,\mathfrak{g}_3] \subset \mathfrak{g}_1 \\
        \\
         (i_1+i_1',i_2+i_2'+1,i_3+i_3'-2) & \text{ if }  \text{ } [\mathfrak{g}_3,\mathfrak{g}_3] \subset \mathfrak{g}_2.
    \end{matrix}\right.  \\
\end{align} % Here $1 \leq \alpha,\beta,i,j \leq 3$ are the indices appearing in the Berezin bracket relations of $\mathfrak{su}^*(4)$ in \eqref{classicalrels}.
 
\end{lemma}

 \begin{proof}
     For any non-zero monomials $p  \in \mathcal{Q}_{i_1 + i_2  + i_3}$ and $q \in \mathcal{Q}_{i_1' + i_2' + i_3'}$ without loss of generality, we may write that $p = Y_1Y_2Y_3$ and $q = Z_1Z_2Z_3,$ where  
        \begin{align*}
         &  Y_1 = x_1^{a_1} \cdots x_u^{a_u},    \text{ } Y_2 = x_{u+1}^{a_{u+1}} \cdots x_{u+v}^{a_{u+v}}, \text{ } Y_3 =  x_{u+v+1}^{a_{u+v+1}} \cdots x_{u+v+w}^{a_{u+v+w}} ; \\
      & Z_1 = x_1^{a_1'} \cdots x_u^{a_u'} ,\text{ } Z_2 = x_{u+1}^{a_{u+1}'} \cdots x_{u+v}^{a_{u+v}'} , \text{ }  Z_3 =  x_{u+v+1}^{a_{u+v+1}'} \cdots x_{u+v+w}^{a_{u+v+w}'}.
    \end{align*} 
    Here $a_1 + \ldots + a_u = i_1, a_{u+1} + \ldots + a_{u+v} = i_2$, $a_{u+v+1} + \ldots + a_{u+v+w} = i_3$ and $a_1' + \ldots + a_u' = i_1', a_{u+1}' + \ldots a_{u+v}' = i_2'$, $a_{u+v+1}' + \ldots + a_{u+v+w}' = i_3'.$ A direct computation shows that 
    \begin{align}
     \{p,q\} = \{Y_1,q\}Y_2Y_3 + Y_1\{Z_2,q\}Z_3 + Y_1 Y_2 \{Y_3,q\}. \label{eq:Leibniz2}
 \end{align}    
 By assumption, the equation $\eqref{eq:Leibniz2}$ becomes 
 \begin{align}
     \{p,q\}    =    & \,    Y_1Y_2 \{Y_3,Z_1\}Z_2 Z_3 +  Y_1Y_2 \{Y_3,Z_2\}Z_1 Z_3 +  Y_1Y_2 \{Y_3,Z_3\}Z_2 Z_1 \, . \label{eq:grading3}
 \end{align}   
 By definition, the grading of $\{p,q\}$ in $\eqref{eq:grading3}$ is equal to the grading of each of the components. From the commutator relations, $\{A_1,B_1\} = 0$. We can then discard this term in $\eqref{eq:grading3}.$ For the rest of the components in $\eqref{eq:grading3}$, we  compute them case by case. Starting from the term  $  Y_1Y_2 \{Y_3,Z_1\}Z_2 Z_3$, a direct computation shows that 
 \begin{align*}
     \mathcal{G} \left(  Y_1Y_2 \{Y_3,Z_1\}Z_2 Z_3\right) = & \, \mathcal{G} \left(\{Y_3,Z_1\}\right) + \mathcal{G} (Z_2) + \mathcal{G} (Z_3) + \mathcal{G} (Y_1) +\mathcal{G} (Y_2)  \\
     = & \, (i_1-1,0,i_3) + ( i_1', 0,0 ) + (0,0,i_3')+ (0,i_2,0)+ (0,i_2',0 ) \\
     = & \, (i_1 +i_1'-1,i_2 +i_2',i_3+i_3') .
 \end{align*} Similarly, we have \begin{align*}
      \mathcal{G} \left(Y_1Y_2 \{Y_3,Z_2\}Z_1 Z_3 \right) = & \, (i_1+i_1',i_2+i_2'-1,i_3+i_3')\\
       \mathcal{G} \left(Y_1Y_2 \{Y_3,Z_3\}Z_2 Z_1\right) = & \, \left\{\begin{matrix}
           ( i_1+i_1'+1,i_2+i_2',i_3+i_3'-2) &  \text{ if } \text{ }[\mathfrak{g}_3,\mathfrak{g}_3] \subset \mathfrak{g}_1 \\
           (i_1+i_1',i_2+i_2'+1,i_3+i_3'-2) & \, \text{ if } \text{ }[\mathfrak{g}_3,\mathfrak{g}_3] \subset \mathfrak{g}_2 .
       \end{matrix}\right.   
 \end{align*} 
 Summing all the terms together, we deduce the grading of $\{p,q\}.$  
 \end{proof}

\begin{thebibliography}{99}
\bibitem[MR2906]{MR2906391} C.~Laurent-Gengoux, A.~Pichereau, and P.~Vanhaecke. \newblock {\em Poisson structures}, volume 347 of {\em Grundlehren der mathematischen {W}issenschaften}. \newblock Springer, Heidelberg, 2013.
\end{thebibliography}
\end{document}
